\begin{figure}[htbp]
\centering
\begin{tikzpicture}[
  x=1cm,y=1cm,font=\small,
  line cap=round,line join=round,
  arrow/.style={-{Stealth[length=1.8mm]},line width=.8pt},
  torsion/.style={arrow,draw=figTeal},
  flip/.style={arrow,draw=figGold,dashed},
  card/.style={draw=figMuted,line width=.5pt},
  note/.style={font=\scriptsize,text=figMuted,align=center}
]
  \path[use as bounding box] (-.1,-.25) rectangle (14.35,7.05);
  \node[anchor=west,font=\small\bfseries] at (0,6.8)
    {(a) Components};
  \node[anchor=west,font=\small\bfseries] at (4.9,6.8)
    {(b) Sheets within cells};
  \node[anchor=west,font=\small\bfseries] at (9.5,6.8)
    {(c) Two loop orders};

  % Reuse the same tube, with orientation and normal coordinates suppressed.
  \fill[figInk!3,rounded corners=6pt] (0,2.5) rectangle (4.2,6.35);
  \node[anchor=west] at (.2,6.05) {$\conf$};
  \pic[scale=.90,transform shape] at (2.0,3.85) {versiontube};
  \node[text=figTeal] at (2.0,6.05) {$\mathcal F$};
  \foreach \x/\y in {.7/2.93} {
    \draw[figLine,fill=white] (\x-.27,\y+.12)--(\x-.27,\y-.12)
      arc[start angle=180,end angle=360,x radius=.27cm,y radius=.1cm]
      --(\x+.27,\y+.12);
    \draw[figLine,fill=white] (\x,\y+.12)
      ellipse[x radius=.27cm,y radius=.1cm];
  }
  \node[note] at (1.4,2.91) {$s\mathcal F$};
  \node[note] at (2.1,2.13) {the same six versions};
  \node at (2.1,1.68) {$S/H\longrightarrow S/G_{12}$};

  % Each version cell contains two local lifts near the reference orbit.
  \foreach \y in {2.15,2.83,3.51,4.19,4.87,5.55} {
    \fill[figInk!4,rounded corners=3pt]
      (4.9,\y-.08) rectangle (7.87,\y+.59);
  }
  \fill[figTeal!13,rounded corners=3pt] (4.9,2.07) rectangle (7.87,2.74);
  \foreach \y/\sheet in {
    2.15/E,2.47/b,2.83/t,3.15/tb,3.51/t^2,3.83/t^2b,
    4.19/u,4.51/ub,4.87/tu,5.19/tub,5.55/t^2u,5.87/t^2ub} {
    \path[card,fill=white]
      (5.05,\y)--(7.25,\y)--(7.75,\y+.20)--(5.55,\y+.20)--cycle;
    \fill[figGold] (5.80,\y+.10) circle[radius=.8pt];
    \node[font=\tiny,anchor=west,inner sep=0pt]
      at (6.05,\y+.10) {$\sheet X$};
  }
  \foreach \y/\version in {2.15/H,2.83/tH,3.51/t^2H,4.19/uH,4.87/tuH,5.55/t^2uH} {
    \draw[figMuted] (7.94,\y)--(8.06,\y)--(8.06,\y+.52)--(7.94,\y+.52);
    \node[anchor=west,font=\scriptsize] at (8.12,\y+.26) {$\version$};
  }
  \draw[arrow,figMuted] (6.35,1.98)--node[right,font=\scriptsize] {quotient} (6.35,1.65);
  \draw[figMuted,fill=figInk!3]
    (6.35,1.15) ellipse[x radius=1.24cm,y radius=.43cm];
  \node[font=\scriptsize] at (7.12,1.16) {$V$};
  \fill[figInk] (6.00,1.15) circle[radius=1.2pt];
  \node[below,font=\scriptsize] at (6.00,1.10) {$[X]$};
  \node at (6.35,.36) {$\mathcal F/G_{12}$};
  \node[note] at (2.1,.95) {$X$ and $gX$ become\\the same point downstairs};
  \draw[arrow,figLine] (3.95,1.15)--(4.85,1.15);

  % Lift a t-loop and a b-loop successively from X.
  % With the left G action upstairs, continuation acts on labels
  % by right multiplication: s -> st or s -> sb.
  \node (e1) at (9.7,5.60) {$X$};
  \node (t1) at (11.65,5.60) {$tX$};
  \node[text=figGold] (tb) at (13.70,5.60) {$tbX$};
  \draw[torsion] (e1)--node[above] {$t$} (t1);
  \draw[flip] (t1)--node[above] {$b$} (tb);

  \node (e2) at (9.7,4.16) {$X$};
  \node (b1) at (11.65,4.16) {$bX$};
  \node[text=figTeal] (bt) at (13.70,4.16) {$btX$};
  \draw[flip] (e2)--node[above] {$b$} (b1);
  \draw[torsion] (b1)--node[above] {$t$} (bt);

  \node at (11.75,3.22) {$btX=t^2bX\ne tbX$};
  \node[note] at (11.75,2.54) {same start, different endpoints};

  % The two loops share a basepoint in the quotient.
  \coordinate (base) at (11.7,1.10);
  \draw[torsion] (base)
    .. controls (10.10,2.40) and (10.05,.02) .. (base);
  \draw[flip] (base)
    .. controls (13.30,.02) and (13.35,2.40) .. (base);
  \fill[figInk] (base) circle[radius=1.2pt];
  \node[below,font=\scriptsize] at (11.7,1.03) {$[X]$};
  \node[text=figTeal] at (10.35,1.52) {$t$};
  \node[text=figGold] at (13.05,1.52) {$b$};
  \node[note] at (11.7,.13) {loops downstairs, lifts above};
\end{tikzpicture}
\caption{The $G_{12}$ tube and version cells of Fig.~\ref{fig:chart-deck},
now viewed through a free covering. Each cell contains two local
sheets over a small quotient neighborhood $V$. Orientations and normal
coordinates are suppressed in the tube only. Loops $t,b$ lift from
$X$ to $tX,bX$. Continuing from $sX$ ends at $stX,sbX$, so order matters.
Loop continuation acts on the right of the sheet label, while deck
transformations act on the left.}
\label{fig:monodromy}
\end{figure}
